\section{Experiments}
\label{sec:experiments}

We evaluate \textbf{BiCAT} in simulation
(Secs.~\ref{sec:benchmark-and-baselines}--\ref{sec:simulation-control}) and on
real robots (Sec.~\ref{sec:real-robot-deployment}). Our experiments
assess both downstream policy performance and, through controlled
ablations, the impact of the key design choices of \textbf{BiCAT}.

\subsection{Simulation Setup}
\label{sec:benchmark-and-baselines}
\label{sec:training-protocol}

\begin{table*}[t]
    \centering
    \caption{
        Success rates on ten RoboTwin 2.0 tasks,
        evaluated over 100 rollouts per task and grouped by inter-arm
        coupling. Avg. denotes the unweighted task average.
        Best and second-best results are bolded
        and underlined, respectively.
    }
    \label{tab:main-tokenizer-results}

    \scriptsize
    \setlength{\tabcolsep}{3.2pt}
    \renewcommand{\arraystretch}{1.15}

    \resizebox{\textwidth}{!}{%
    \begin{tabular}{@{}l|ccccccccccc@{}}
        \toprule
        & \multicolumn{2}{c}{Single-arm}
        & \multicolumn{4}{c}{Loosely coupled}
        & \multicolumn{4}{c}{Tightly coupled} & \\
        \cmidrule(lr){2-3}\cmidrule(lr){4-7}\cmidrule(lr){8-11}

        \shortstack[c]{Tokenizer}
        &
        \shortstack[c]{Shake bottle\\horizontally}
        &
        \shortstack[c]{Place container\\plate}
        &
        \shortstack[c]{Stack blocks\\two}
        &
        \shortstack[c]{Dump bin\\bigbin}
        &
        \shortstack[c]{Stack bowls\\three}
        &
        \shortstack[c]{Pick dual\\bottles}
        &
        \shortstack[c]{Hand over\\mic}
        &
        \shortstack[c]{Place bread\\skillet}
        &
        \shortstack[c]{Place dual\\shoes}
        &
        \shortstack[c]{Put bottles\\dustbin}
        &
        \shortstack[c]{Avg.}
        \\
        
        \midrule

        Bin~\cite{kim2024openvla}
        & 0.67 & 0.30 & 0.00 & 0.00 & 0.00 & 0.00 & 0.05 & 0.00 & 0.00 & 0.00 & 0.102 \\
        FAST~\cite{pertsch2025fast}
        & 0.70 & 0.25 & 0.15 & 0.27 & 0.09 & 0.11 & 0.52 & 0.13 & 0.01 & 0.00 & 0.223 \\
        QueST~\cite{mete2024quest}
        & 0.96 & \textbf{0.92} & 0.43 & 0.69 & 0.55 & 0.12 & \underline{0.90} & 0.10 & 0.04 & 0.20 & 0.491 \\
        $\mathrm{OAT}^{\mathrm{sing}}$~\cite{liu2026oat_rss}
        & 0.93 & \underline{0.90} & 0.47 & 0.78 & 0.50 & 0.03 & \textbf{0.94} & 0.08 & 0.03 & 0.21 & 0.487 \\
        $\mathrm{OAT}^{\mathrm{pow2}}$~\cite{liu2026oat_rss}
        & 0.93 & 0.85 & \underline{0.49} & 0.70 & 0.54 & 0.10 & 0.77 & 0.06 & 0.05 & 0.20 & 0.469 \\
        $\mathcal{M}^{2}$Tok~\cite{xu2026m2tok}
        & 0.87 & 0.69 & 0.18 & 0.53 & 0.43 & 0.12 & 0.83 & 0.14 & 0.05 & 0.21 & 0.405 \\
        NAC~\cite{jawaid2026nac}
        & \underline{0.99} & \underline{0.90} & 0.47 & \textbf{0.87} & \underline{0.57} & 0.15 & 0.74 & 0.18 & 0.04 & 0.26 & \underline{0.517} \\
        ActionCodec~\cite{dong2026actioncodec}
        & 0.98 & 0.83 & 0.34 & 0.60 & 0.54 & \underline{0.23} & 0.72 & \underline{0.25} & \underline{0.06} & \underline{0.37} & 0.492 \\
        \oursrow\textbf{BiCAT (Ours)}
        & \textbf{1.00} & \textbf{0.92} & \textbf{0.66} & \underline{0.85} & \textbf{0.81} & \textbf{0.44} & 0.88 & \textbf{0.47} & \textbf{0.41} & \textbf{0.49} & \textbf{0.693} \\
        \bottomrule
    \end{tabular}%
    }

    \vspace{-2mm}
\end{table*}

\begin{figure*}[t]
    \centering
    \includegraphics[width=\textwidth]{assets/why_bicat_success_jump.pdf}
    \caption{\textbf{What tracks closed-loop success.} Filled markers: nine
    tokenizers; hollow grey markers: seven BiCAT ablations
    (Table~\ref{tab:ablation}); y-axis: mean RoboTwin success. $\rho_{16}$ and
    $\rho_{9}$ are Spearman correlations over all 16 points and over the nine
    tokenizers.
    (a)~Reconstruction MSE of held-out chunks in the original action space
    ($\times10^{-3}$; $\mathcal{M}^{2}$Tok, 2.25, is clipped).
    (b)~Policy token accuracy (\%).
    (c)~Ratio of action-token NLL with shuffled versus matched images; higher
    means stronger reliance on vision.
    (d)~Change of the decoded action when $1/16$ of the tokens are replaced
    ($\times10^{-3}$, log scale).
    (e)~Disagreement between chunks predicted 8 steps apart over their 24
    overlapping steps ($\times10^{-3}$).
    (f)~Mean absolute change of the arm joints at chunk boundaries
    ($\times10^{-3}$\,rad). Lower is better in (d)--(f).}
    \label{fig:why-bicat-success}
\end{figure*}



\paragraph{Benchmark and protocol.}
RoboTwin 2.0~\cite{chen2025robotwin2} comprises 50 bimanual manipulation
tasks. Tokenizers and policies are trained on all 2,500 demonstrations
from all 50 tasks. Each learned tokenizer is pretrained with its own
objective and then frozen, and a PaliGemma~2-3B
backbone~\cite{steiner2024paligemma2} is fully fine-tuned to predict its
action tokens from a language instruction, the discretized robot state, and
three RGB views. Policies predict 32-step chunks and execute the first 16
before replanning. As many of these tasks need only one arm, we
evaluate on ten representative tasks, selected before training
to cover three levels of inter-arm coupling, inspired by the coordination
taxonomy of Sa~\etal~\cite{sa2026vlareview}: two \emph{single-arm} tasks,
four \emph{loosely coupled} tasks where both arms act largely
independently, and four \emph{tightly coupled} tasks where each arm's
motion depends on its partner (Appendix~\ref{app:robotwin-tasks}). We
report success over 100 rollouts per task and the ten-task mean
(settings in Appendix~\ref{app:baselines}).

\paragraph{Baselines.}
We compare \textbf{BiCAT} with Bin~\cite{kim2024openvla},
FAST~\cite{pertsch2025fast}, QueST~\cite{mete2024quest},
OAT~\cite{liu2026oat_rss} with token-wise ($\mathrm{OAT}^{\mathrm{sing}}$)
and five-block ($\mathrm{OAT}^{\mathrm{pow2}}$) prediction,
$\mathcal{M}^{2}$Tok~\cite{xu2026m2tok}, NAC~\cite{jawaid2026nac}, and
ActionCodec~\cite{dong2026actioncodec}; Table~\ref{tab:efficiency} lists
their token budgets, vocabularies, sizes, prediction schedules, and
latencies. To keep the comparison controlled, all learned tokenizers use a
vocabulary of 1920 or 2048 codes, 16 tokens per chunk except NAC (24,
following its multi-scale design), and between 7.4M and 11.9M parameters,
with \textbf{BiCAT} the smallest. All tokenizers are fitted or trained on
the RoboTwin training data; FAST is refitted rather than taken from its
released universal tokenizer, and $\mathcal{M}^{2}$Tok uses its released
implementation (Appendix~\ref{app:tokenizer-configurations}). All methods
share the data, backbone,
observations, and optimization settings and differ only in the tokenizer
and its prediction schedule.



\subsection{Simulation Results and Analysis}
\label{sec:simulation-control}

We use RoboTwin 2.0 to answer four research questions.
\textbf{(Q1)}~Does building bimanual structure into the tokenizer improve
bimanual VLA control over existing tokenizers?
\textbf{(Q2)}~Where does the improvement come from?
\textbf{(Q3)}~How sensitive is \textbf{BiCAT} to its vocabulary size, and
where does it fall short?
\textbf{(Q4)}~What does two-block prediction cost at inference time?

\begin{table}[!t]
    \centering
    \caption{Token budgets, prediction schedules, tokenizer size, and
    per-chunk latency on one H200 GPU. $K$: vocabulary size;
    $[b]\!\times\!n$: $n$ blocks of $b$ parallel tokens; Param.: tokenizer
    parameters (M), excluding NAC's training-only discriminator; Rec.:
    separately timed encode--decode round trip; Policy: mean token-generation
    time per chunk.}
    \label{tab:efficiency}
    \footnotesize
    \setlength{\tabcolsep}{3pt}
    \renewcommand{\arraystretch}{1.08}
    \resizebox{\columnwidth}{!}{%
    \begin{tabular}{@{}lcccccccc@{}}
        \toprule
        & & & & & \multicolumn{3}{c}{Tokenizer (ms)} & Policy \\
        \cmidrule(lr){6-8}
        Tokenizer & \#Tok. & $K$ & Schedule & Param. & Encode & Decode & Rec. & (ms) \\
        \midrule
        Bin & 448 & 256 & $[14]\!\times\!32$ & -- & 0.07 & 0.06 & 0.13 & 451.7 \\
        FAST & $N$ & 2048 & $[1]\!\times\!N$ & -- & 0.37 & 0.22 & 0.65 & 533.6 \\
        QueST & 16 & 1920 & $[1]\!\times\!16$ & 11.61 & 1.74 & 1.58 & 3.44 & 227.7 \\
        $\mathrm{OAT}^{\mathrm{sing}}$ & 16 & 1920 & $[1]\!\times\!16$ & 11.04 & 2.95 & 1.98 & 4.94 & 229.1 \\
        $\mathrm{OAT}^{\mathrm{pow2}}$ & 16 & 1920 & $[1,1,2,4,8]$ & 11.04 & 2.90 & 1.98 & 4.88 & 79.6 \\
        $\mathcal{M}^{2}$Tok & 16 & 2048 & $[8,8]$ & 11.27 & 2.05 & 1.98 & 3.92 & 40.6 \\
        NAC & 24 & 2048 & $[8,16]$ & 8.48 & 1.42 & 2.41 & 3.92 & 40.8 \\
        ActionCodec & 16 & 2048 & $[16]$ & 11.87 & 1.56 & 1.32 & 2.87 & 26.8 \\
        \midrule
        \textbf{BiCAT} & 16 & 1920 & $[8,8]$ & 7.40 & 1.74 & 1.64 & 3.45 & 40.5 \\
        \bottomrule
    \end{tabular}%
    }
    \vspace{-2mm}
\end{table}


\paragraph{Q1: How does BiCAT compare with baselines?}
As shown in Table~\ref{tab:main-tokenizer-results}, \textbf{BiCAT} reaches a mean
success rate of 69.3\%, compared with 51.7\% for NAC, the strongest
baseline. The gain of 17.6 percentage points
is accompanied by the best or tied-best result on eight of the ten tasks,
and it grows with inter-arm coupling: over the strongest baseline in each
group, \textbf{BiCAT} gains 1.5 points on the near-saturated single-arm
tasks, 17.5 on loosely coupled tasks, and 21.3 on tightly coupled tasks.
Since the learned tokenizers have comparable token budgets
(Sec.~\ref{sec:benchmark-and-baselines}), the gain is not explained by a shorter action sequence.
$\mathcal{M}^{2}$Tok, which also tokenizes each arm separately but with a
different architecture and multi-codebook quantization, reaches 40.5\%; we
isolate the role of cross-arm interaction with controlled ablations in Q2.




\paragraph{Q2: What explains the gains?}
Figure~\ref{fig:why-bicat-success} relates mean success to six properties of
each tokenizer and its trained policy. Consistent with
ActionCodec~\cite{dong2026actioncodec}, reconstruction fidelity does not
predict control ($\rho_{16}=-0.13$): Bin and FAST reconstruct actions most
accurately yet perform worst. Token accuracy is also uninformative
($\rho_{9}=-0.42$), as it is not comparable across vocabularies. Success
instead tracks reliance on visual input ($\rho_{16}=0.72$), consistency
across replanning ($\rho_{16}=-0.80$), and smoothness at chunk boundaries
($\rho_{16}=-0.78$), and less closely tolerance to token errors
($\rho_{16}=-0.51$); \textbf{BiCAT} has the smallest boundary jump among the
nine tokenizers. As eight of the 16 points are \textbf{BiCAT} or its
variants, we read these correlations as descriptive rather than causal.

\paragraph{Q2: Which components contribute?}
Table~\ref{tab:ablation} ablates the design principle and the prediction
schedule. \emph{Raw coordinates} sets $\mathbf{s}=\mathbf{1}$, so the two
arms share features for identical joint values rather than for mirrored
motions, and a motion such as reaching inward is encoded differently for
each arm; \emph{Arm-specific modules} gives each arm its own modules and
normalization bounds. These variants lower success to 58.7\% and 56.9\%;
sharing
aligns the latents of the two arms (left--right MMD of 0.088 for
\textbf{BiCAT}, versus 0.193 and 0.286; Appendix~\ref{app:arm-latents}). Removing the
pre-quantization cross-attention (\emph{Decoder-only coupling}), joint
decoding (\emph{Encoder-only coupling}), or both (\emph{Independent arms}, a
controlled per-arm variant) lowers success to
51.7\%, 56.0\%, and 50.0\%. The two interactions are complementary: each
adds far more when the other is present (17.6 versus 6.0 points for
cross-attention, 13.3 versus 1.7 for joint decoding), and removing them
weakens error tolerance and visual grounding. Predicting both arms in one
block (\emph{Single-block prediction}) lowers success to 51.1\%, so neither
partner-aware codes without the arm-wise schedule nor the schedule without
partner-aware codes exceeds NAC. Reversing the block order
(\emph{Right-first prediction}) lowers success to 54.2\%, so the order also
matters; we use left-first throughout.


% Embodiment groups run left to right; one tokenizer per row.




\begin{table}[t]
    \centering
    % DRAFT: Err. tol. and Vis. ground. come from the bicat_why analysis; verify before submission.
    \caption{Ablations on the ten RoboTwin 2.0 tasks. Columns mark the retained
    components: coordinate canonicalization, cross-arm sharing of modules and
    normalization, pre-quantization cross-attention, and joint decoding; Order
    is the prediction schedule (L$\to$R: left-arm block first; R$\to$L:
    right-arm block first; 1 blk: both arms in one block). SR is mean success (\%). Err.\ tol.\ is the change of the decoded
    action when $1/16$ of the tokens are replaced ($\times10^{-3}$, lower is
    better), and Vis.\ is the visual-grounding ratio of
    Fig.~\ref{fig:why-bicat-success}c (higher is better). Per-task results are
    in Appendix~\ref{app:ablation-per-task}.}
    \label{tab:ablation}
    \footnotesize
    \setlength{\tabcolsep}{2.6pt}
    \renewcommand{\arraystretch}{1.1}
    \newcommand{\yes}{$\checkmark$}
    \newcommand{\no}{--}
    \resizebox{\columnwidth}{!}{%
    \begin{tabular}{@{}lccccc|ccc@{}}
        \toprule
        Variant & Can. & Share & CA & Joint & Order & SR & Err.\ tol. & Vis. \\
        \midrule
        Raw coordinates & \no & \yes & \yes & \yes & L$\to$R & 58.7 & 7.7  & 6.1 \\
        Arm-specific modules & \yes & \no & \yes & \yes & L$\to$R & 56.9 & 7.5  & 4.8 \\
        \midrule
        Decoder-only coupling & \yes & \yes & \no & \yes & L$\to$R & 51.7 & 10.0 & 4.3 \\
        Encoder-only coupling & \yes & \yes & \yes & \no & L$\to$R & 56.0 & 22.9 & 3.8 \\
        Independent arms & \yes & \yes & \no & \no & L$\to$R & 50.0 & 24.1 & 4.4 \\
        \midrule
        Single-block prediction & \yes & \yes & \yes & \yes & 1 blk & 51.1 & 8.1  & 7.8 \\
        Right-first prediction   & \yes & \yes & \yes & \yes & R$\to$L & 54.2 & 8.1 & 6.6 \\
        \midrule
        \textbf{BiCAT} & \yes & \yes & \yes & \yes & L$\to$R & \textbf{69.3} & \textbf{7.3} & \textbf{7.9} \\
        \bottomrule
    \end{tabular}%
    }
    \vspace{-2mm}
\end{table}

\begin{table}[t]
    \centering
    \caption{FSQ grid ablation. Mean success rate over the ten RoboTwin 2.0
    tasks for different FSQ levels and their induced vocabulary sizes.}
    \label{tab:FSQ}
    \footnotesize
    \setlength{\tabcolsep}{3.5pt}
    \renewcommand{\arraystretch}{1.1}
    \resizebox{0.95\columnwidth}{!}{%
    \begin{tabular}{@{}l|ccccc@{}}
        \toprule
        FSQ Levels
        & [8,6,5]
        & [8,8,8]
        & [8,5,5,5]
        & [8,8,6,5]
        & [7,5,5,5,5] \\
        Induced $|\mathcal{V}|$
        & 240 & 512 & 1000 & 1920 & 4375 \\
        \midrule
        Mean SR
        & 0.470
        & 0.617
        & 0.631
        & 0.693
        & 0.589 \\
        \bottomrule
    \end{tabular}%
    }
\end{table}

\begin{table*}[t]
    \centering
    \caption{Real-world task success and inference latency, grouped by embodiment.
    Success is reported as successful trials over evaluated trials.
    Overall pools trials across all three embodiments, and latency reports
    the YAM measurements.}
    \label{tab:real-world-results}
    \label{tab:three-task-success-counts}
    \label{tab:unitree-task-success-counts}
    \label{tab:so101-task-success-counts}
    \footnotesize
    \setlength{\tabcolsep}{3.5pt}
    \renewcommand{\arraystretch}{1.15}
    \begin{tabularx}{\textwidth}{@{}l*{7}{>{\centering\arraybackslash}X}>{\centering\arraybackslash}p{2.0cm}@{}}
        \toprule
        & \textbf{SO-101}
        & \multicolumn{2}{c}{\textbf{Unitree G1}}
        & \multicolumn{3}{c}{\textbf{YAM}}
        & \textbf{Overall} & \textbf{Latency} \\
        \cmidrule(lr){2-2}
        \cmidrule(lr){3-4}
        \cmidrule(lr){5-7}
        \cmidrule(lr){8-8}
        \cmidrule(l){9-9}
        Tokenizer & Block & Bottle & Bowl & Pen & Box & Bottle & All tasks & (ms) \\
        \midrule
        FAST~\cite{pertsch2025fast}
        & 3/25 & 6/25 & 7/25 & 0/25 & 1/25 & 8/25 & 25/150 & $434.1 \pm 103.6$ \\
        $\mathrm{OAT}^{\mathrm{sing}}$~\cite{liu2026oat_rss}
        & 19/25 & 13/25 & 15/25 & 5/25 & 7/25 & 15/25 & 74/150 & $251.6 \pm 5.1$ \\
        \oursrow\textbf{BiCAT (Ours)}
        & \textbf{25/25} & \textbf{24/25} & \textbf{21/25} & \textbf{8/25} & \textbf{14/25} & \textbf{21/25} & \textbf{113/150} & \textbf{91.7 $\pm$ 1.3} \\
        \bottomrule
    \end{tabularx}
\end{table*}


\begin{figure*}[t]
    \centering
    \includegraphics[width=\textwidth]{assets/real_rollouts.pdf}
    \caption{\textbf{BiCAT rollouts on three real bimanual embodiments.}
    Each row shows four keyframes of one task with a short task description.}
    \label{fig:real-rollouts}
\end{figure*}

\paragraph{Q3: How does the vocabulary affect control?}
Table~\ref{tab:FSQ} varies the FSQ levels of \textbf{BiCAT} (per-task results
in Appendix~\ref{app:fsq-per-task}). We choose the grids following
FSQ~\cite{mentzer2024fsq}: all levels satisfy $L_j\geq5$, and $[8,6,5]$,
$[8,5,5,5]$, and $[7,5,5,5,5]$ are its recommended configurations for
vocabularies of about $2^{8}$, $2^{10}$, and $2^{12}$. Success rises from
47.0\% with 240 codes to 69.3\% with 1920 codes and falls to 58.9\% with
4375 codes, so additional quantization capacity does not monotonically
improve control, although changing the levels also changes the latent
dimensionality. All other experiments use $[8,8,6,5]$, as do QueST and OAT.

\paragraph{Q3: Where does BiCAT underperform?}
The improvement is not uniform. \textbf{BiCAT} reaches 88\% on
\textsc{Handover Mic}, below $\mathrm{OAT}^{\mathrm{sing}}$ (94\%) and QueST
(90\%), and 85\% on \textsc{Dump Bin Bigbin}, below NAC (87\%). On
\textsc{Handover Mic}, four tokenizer ablations with the same block order
reach 92--96\% (Appendix~\ref{app:ablation-per-task}), so this gap is not
inherent to arm-wise tokens or two-block prediction.

\paragraph{Q4: What does two-block prediction cost?}
Table~\ref{tab:efficiency} reports per-chunk latency on one H200 GPU with
batch size 1. Encoding and decoding take under 5\,ms for every tokenizer,
so latency is set by the prediction schedule. With two blocks, \textbf{BiCAT} takes 40.5\,ms, comparable to
$\mathcal{M}^{2}$Tok and NAC, which also predict two blocks, and
$5.6\times$ to over $11\times$ faster than tokenizers decoded token by token.
Only ActionCodec, which predicts a single block, is faster (26.8\,ms), but
\emph{Single-block prediction} costs \textbf{BiCAT} 18.2 points of success,
so the second block is worth its extra decoding step.


\subsection{Deployment on Real Robots}
\label{sec:real-robot-deployment}

\paragraph{Platforms and tasks.}
We evaluate on three bimanual embodiments: Bimanual SO-101, the Unitree G1
humanoid, and Bimanual YAM. The suite contains six tasks: Block on SO-101,
Bottle and Bowl on G1, and Pen, Box, and Bottle on YAM.
Fig.~\ref{fig:real-rollouts} shows each task, and
Appendix~\ref{app:real-world} gives task statistics.

\paragraph{Data and training.}
We collect 100 teleoperated demonstrations per task (105 for Pen). For each
embodiment, one tokenizer is trained on all of its demonstrations and
frozen, and one VLA policy with a Molmo2-ER
backbone~\cite{fang2026molmoact2} is trained on its tokens for all tasks of
that embodiment. Within an embodiment, all methods use exactly the same
data and training configuration, including training steps, learning rate,
warmup, and batch size (Appendix~\ref{app:real-world-training}), and the
final checkpoint is evaluated.

\paragraph{Baselines and evaluation.}
We compare \textbf{BiCAT} with one representative tokenizer from each of the
two main families: FAST~\cite{pertsch2025fast}, the widely adopted
compression-based tokenizer, and $\mathrm{OAT}^{\mathrm{sing}}$~\cite{liu2026oat_rss},
a learned tokenizer with ordered tokens and token-wise prediction. All
policies run on one NVIDIA RTX 5090 GPU and are evaluated for 25 trials per
task, 150 per method. A trial succeeds only if every stage of the task is
completed, and the initial object placements are identical across methods.
Inference latency is measured on YAM, from receiving an observation to
outputting an action chunk, over all inference calls.
The robot waits during inference, so latency lengthens an episode but does
not change the observations on which chunks are predicted
(Appendix~\ref{app:real-world-eval}).

\paragraph{Results.}
As shown in Table~\ref{tab:real-world-results}, \textbf{BiCAT} succeeds in 113
of 150 trials, compared with 74 for $\mathrm{OAT}^{\mathrm{sing}}$ and 25 for
FAST, leads on all six tasks, and has the lowest inference latency, $4.7\times$
lower than FAST. Consistent with simulation, the advantage persists across physical bimanual embodiments.
